% ============================================================
% 66_index.tex — OPTIONALES Stichwortverzeichnis (alphabetischer Index)
% ============================================================
% Opt-in: in preamble.tex einkommentieren, wenn die ZSF ein
% Stichwortverzeichnis bekommen soll. Das Modul lädt makeidx selbst.
% Weitere Voraussetzungen: \ZSF@refExtractChap + \ZSFkeyword
% (50_typography_semantics), \ZSFspace* (30_layout_spacing),
% \ZSFfontBoxTitle (50). makeindex-Style: styles/zsfindex.ist
% (Makefile übergibt ihn via INDEXSTYLE + $makeindex an latexmk).
%
% Herkunft: konsolidiert aus zsf-linalg-fs2026/styles/65_index.tex.
%
% Index-Einträge verweisen auf ABSCHNITTSNUMMERN ("6.2" / "6.2.1") und
% standardmässig zusätzlich auf die physische Seite. Die Abschnittsnummer
% bleibt dabei der primäre Wegweiser und trägt die Kapitel-Palettenfarbe.
%
% Mechanik: Die Abschnittsnummer wird per \edef SOFORT eingefroren —
% beim Shipout wäre \thesubsection längst weitergezählt. Die Seitenzahl
% dagegen MUSS verzögert bleiben (\protected@write + \noexpand\thepage):
% \thepage ist an der Schreibstelle noch die zuletzt ausgegebene Seite,
% nicht die, auf der das Material landet. Unter multicols* sammelt eine
% Seite vier Spalten, wodurch ein sofort expandiertes \thepage
% systematisch zu klein ausfällt. makeindex sortiert mit
% styles/zsfindex.ist (page_compositor ".") und merged identische
% Duplikate aus tabularx/tblr-Mehrfach-Trials automatisch.
%
% Verwendung am Dokumentende: eigenes Kapitel mit
%   \StartFrontChapter{Stichwortverzeichnis}
%   \printindex
%
% Öffentliche Makros (Verwendung in Kapiteln, siehe rules/55_index.md):
%   \ZSFkeyword{X}              — fett + automatischer Indexeintrag
%   \ZSFkeyword*{X}             — nur fett, KEIN Indexeintrag
%   \ZSFkeyword[sort@Anzeige]{X}— Eintrag-Override (Umlaut-Sortkey,
%                                 abweichende Indexform)
%   \ZSFindex[sortkey]{Begriff} — unsichtbarer Eintrag
%   \ZSFindexsee{Synonym}{Ziel} — Verweis "Synonym, siehe Ziel 6.2"
%                                 (mit farbcodierter Zielnummer, Ein-Hop)

\RequirePackage{makeidx}
\makeindex
\renewcommand{\seename}{siehe}
\newif\ifZSFIndexShowPageNumber
\ZSFIndexShowPageNumbertrue

% ------------------------------------------------------------
% Writer: Nummer sofort, Seitenzahl verzögert
% ------------------------------------------------------------
\makeatletter
% Tiefste aktuelle Abschnittsnummer einfrieren; vor der ersten
% \SubsectionBar eines Kapitels Fallback auf die Kapitelnummer.
\newcommand{\ZSF@idxnumnow}{%
  \ifnum\value{subsubsection}>0 \thesubsubsection\else
    \ifnum\value{subsection}>0 \thesubsection\else\thesection\fi
  \fi}
% #1 = Eintrag (ggf. sort@display), #2 = Encap (ZSFidxnum | see{...}).
% Im Seitenzahl-Modus wird ein zusammengesetzter Locator Abschnitt.Seite
% geschrieben. makeindex kann ihn mit page_compositor "." sortieren und
% verdichten; die Ausgabe trennt beide Teile erst am Schluss wieder.
% Der Eintragstext liegt in Token-Registern, damit \edef nur den Locator
% zusammensetzt und Makros im Begriff nicht vorzeitig expandiert.
\newcommand{\ZSF@writeidx}[2]{%
  \@ifundefined{@indexfile}{}{%
    \begingroup
      \toks0={#1}%
      \toks1={#2}%
      \toks2={\protected@write\@indexfile{}}%
      \edef\ZSF@num{\ZSF@idxnumnow}%
      \ifZSFIndexShowPageNumber
        \edef\ZSF@line{%
          \the\toks2{\string\indexentry{\the\toks0|\the\toks1}%
            {\ZSF@num.\noexpand\thepage}}}%
      \else
        \edef\ZSF@line{%
          \the\toks2{\string\indexentry{\the\toks0|\the\toks1}{\ZSF@num}}}%
      \fi
      \ZSF@line
    \endgroup}}
% «siehe»-Verweise bleiben reine Abschnitts-Wegweiser und erhalten nie
% eine Seitenzahl.
\newcommand{\ZSF@writeidxsee}[2]{%
  \@ifundefined{@indexfile}{}{%
    \begingroup
      \toks0={#1}%
      \toks1={#2}%
      \toks2={\protected@write\@indexfile{}}%
      \edef\ZSF@num{\ZSF@idxnumnow}%
      \edef\ZSF@line{%
        \the\toks2{\string\indexentry{\the\toks0|see{\the\toks1}}{\ZSF@num}}}%
      \ZSF@line
    \endgroup}}

% ------------------------------------------------------------
% Öffentliche Eintrags-Makros
% ------------------------------------------------------------
% (bleiben im \makeatletter-Block: Bodies referenzieren \ZSF@writeidx)
\NewDocumentCommand{\ZSFindex}{o m}{%
  \IfNoValueTF{#1}%
    {\ZSF@writeidx{#2}{ZSFidxnum}}%
    {\ZSF@writeidx{#1@#2}{ZSFidxnum}}}
\NewDocumentCommand{\ZSFindexsee}{o m m}{%
  \IfNoValueTF{#1}%
    {\ZSF@writeidxsee{#2}{#3}}%
    {\ZSF@writeidxsee{#1@#2}{#3}}}

% \ZSFkeyword (Basisdefinition in 50_typography_semantics) indexiert jetzt
% automatisch mit. Die SIGNATUR ist identisch zur Basis — dieses Modul ergänzt
% nur das Schreiben, es ändert die Aufrufform nicht. Das Rendering läuft
% weiterhin über \ZSFkeywordStyle, damit ein Fork die Darstellung an einer
% Stelle ändern kann, unabhängig davon, ob der Index aktiv ist.
\RenewDocumentCommand{\ZSFkeyword}{s o m}{%
  \ZSFkeywordStyle{#3}%
  \IfBooleanF{#1}{%
    \IfNoValueTF{#2}{\ZSFindex{#3}}{\ZSFindex{#2}}}}
\makeatother

% ------------------------------------------------------------
% Rendering der Eintragsnummern (Encap)
% ------------------------------------------------------------
% makeindex emittiert \ZSFidxnum{6.2} bzw. im Seitenzahl-Modus
% \ZSFidxnum{6.2.14}; Kapitelziffer-Extraktion wie bei
% \ZSFref (\ZSF@refExtractChap aus 50_typography_semantics).
\makeatletter
\newcommand{\ZSF@idxrendersection}[1]{%
  \begingroup
    \edef\ZSF@idxTargetChap{\expandafter\ZSF@refExtractChap#1.\@nil}%
    \ZSFSetPaletteColorNameFromNumber{\ZSF@idxTargetChap}{\ZSF@idxTargetColor}%
    {\upshape\bfseries\color{\ZSF@idxTargetColor}#1}%
  \endgroup}
\ExplSyntaxOn
\seq_new:N \l__zsf_index_locator_seq
\tl_new:N \l__zsf_index_section_tl
\tl_new:N \l__zsf_index_page_tl
\tl_new:N \l__zsf_index_pending_section_tl
\tl_new:N \l__zsf_index_pending_page_tl
\bool_new:N \l__zsf_index_locator_multiple_bool
\cs_new_protected:Npn \ZSFIndexLocatorListStart
  {\tl_clear:N \l__zsf_index_pending_section_tl \tl_clear:N \l__zsf_index_pending_page_tl \bool_set_false:N \l__zsf_index_locator_multiple_bool}
\cs_new_protected:Npn \ZSFIndexLocatorListMultiple
  {
    \bool_set_true:N \l__zsf_index_locator_multiple_bool
    \tl_if_empty:NF \l__zsf_index_pending_section_tl
      {\exp_args:NV \ZSF@idxrendersection \l__zsf_index_pending_section_tl}
    \tl_clear:N \l__zsf_index_pending_section_tl
    \tl_clear:N \l__zsf_index_pending_page_tl
  }
\cs_new_protected:Npn \ZSFIndexLocatorListFinish
  {
    \bool_if:NF \l__zsf_index_locator_multiple_bool
      {\tl_if_empty:NF \l__zsf_index_pending_section_tl
        {\mbox{\exp_args:NV \ZSF@idxrendersection \l__zsf_index_pending_section_tl\ZSFindexLocatorSep\textperiodcentered\ZSFindexLocatorSep(S.~\tl_use:N \l__zsf_index_pending_page_tl)}}}
    \tl_clear:N \l__zsf_index_pending_section_tl
    \tl_clear:N \l__zsf_index_pending_page_tl
  }
\cs_new_protected:Npn \ZSFIndexRenderLocator #1
  {
    \seq_set_split:Nnn \l__zsf_index_locator_seq {.} {#1}
    \seq_pop_right:NN \l__zsf_index_locator_seq \l__zsf_index_page_tl
    \tl_set:Nx \l__zsf_index_section_tl
      {\seq_use:Nn \l__zsf_index_locator_seq {.}}
    \bool_if:NTF \l__zsf_index_locator_multiple_bool
      {\exp_args:NV \ZSF@idxrendersection \l__zsf_index_section_tl}
      {\tl_set:Nx \l__zsf_index_pending_section_tl {\tl_use:N \l__zsf_index_section_tl}\tl_set:Nx \l__zsf_index_pending_page_tl {\tl_use:N \l__zsf_index_page_tl}}
  }
\ExplSyntaxOff
\newcommand{\ZSFidxnum}[1]{%
  \ifZSFIndexShowPageNumber
    \ZSFIndexRenderLocator{#1}%
  \else
    \ZSF@idxrendersection{#1}%
  \fi}
\newcommand{\ZSFidxsee}[1]{\ZSF@idxrendersection{#1}}
\makeatother

% see-Verweise zeigen zusätzlich die (farbcodierte) Abschnittsnummer des
% Ziels an, damit ein Synonym-Lookup in EINEM Schritt landet — ohne diese
% Zeile griffe makeidx' Default \see, der die Nummer (#2) verwirft.
% makeindex liefert \see{Ziel}{Nummer}; #2 ist die am kanonischen Ort
% eingefrorene Zielnummer (Konvention: \ZSFindexsee am Zielort platzieren).
\renewcommand*{\see}[2]{\emph{\seename} #1\enspace\ZSFidxsee{#2}}

% ------------------------------------------------------------
% theindex-Umgebung (Override) + Buchstaben-Heading
% ------------------------------------------------------------
% article-Default macht \twocolumn — illegal in multicols*. Stattdessen
% einfache hängende Liste im laufenden 4-Spalten-Fluss.
\makeatletter
\newcommand{\ZSF@idxitem}{\par\hangindent\ZSFindexItemIndent\relax}
\newcommand{\ZSF@idxsubitem}{\par\hangindent\ZSFindexSubitemIndent\hspace*{\ZSFindexSubitemOffset}}
\newcommand{\ZSF@idxsubsubitem}{\par\hangindent\ZSFindexSubsubitemIndent\hspace*{\ZSFindexSubsubitemOffset}}
\newcommand{\ZSF@idxspace}{\par\addvspace{\ZSFspaceM}}
\renewenvironment{theindex}{%
  \ZSFTitleBarLockOff % Register ist Inhalt: Balken-Sperre endet hier
  \begingroup
  \RaggedRight
  \setlength{\parindent}{0pt}%
  \setlength{\parskip}{0pt plus \ZSFindexParStretch}%
  \let\item\ZSF@idxitem
  \let\subitem\ZSF@idxsubitem
  \let\subsubitem\ZSF@idxsubsubitem
  \let\indexspace\ZSF@idxspace
}{\par\endgroup}
\makeatother
\newcommand{\ZSFIndexLetter}[1]{%
  \par\addvspace{\ZSFspaceM}%
  {\ZSFfontBoxTitle #1}\par\nobreak\addvspace{\ZSFspaceXS}}
